%-------------------------
% Resume in Latex
% Based off of: https://github.com/jakegut/resume (Jake Gutierrez, MIT)
%------------------------

\documentclass[letterpaper,11pt]{article}

\usepackage{latexsym}
\usepackage[empty]{fullpage}
\usepackage{titlesec}
\usepackage{marvosym}
\usepackage[usenames,dvipsnames]{color}
\usepackage{verbatim}
\usepackage{enumitem}
\usepackage[hidelinks]{hyperref}
\usepackage{fancyhdr}
\usepackage[french]{babel}
\frenchsetup{StandardLists=true}
\usepackage{tabularx}
\usepackage[T1]{fontenc}
\usepackage[utf8]{inputenc}
\input{glyphtounicode}

\pagestyle{fancy}
\fancyhf{}
\fancyfoot{}
\renewcommand{\headrulewidth}{0pt}
\renewcommand{\footrulewidth}{0pt}

% A4 page, adjusted margins
\usepackage[a4paper,top=0.45in,bottom=0.45in,left=0.5in,right=0.5in]{geometry}

\urlstyle{same}

\raggedbottom
\raggedright
\setlength{\tabcolsep}{0in}

% Sections formatting
\titleformat{\section}{
  \vspace{-4pt}\scshape\raggedright\large
}{}{0em}{}[\color{black}\titlerule \vspace{-5pt}]

% Ensure that generate pdf is machine readable/ATS parsable
\pdfgentounicode=1

%-------------------------
% Custom commands
\newcommand{\resumeItem}[1]{
  \item\small{
    {#1 \vspace{-2pt}}
  }
}

\newcommand{\resumeSubheading}[4]{
  \vspace{-2pt}\item
    \begin{tabular*}{0.97\textwidth}[t]{l@{\extracolsep{\fill}}r}
      \textbf{#1} & #2 \\
      \textit{\small#3} & \textit{\small #4} \\
    \end{tabular*}\vspace{-7pt}
}

\newcommand{\resumeProjectHeading}[2]{
    \item
    \begin{tabular*}{0.97\textwidth}{l@{\extracolsep{\fill}}r}
      \small#1 & #2 \\
    \end{tabular*}\vspace{-7pt}
}

\renewcommand\labelitemii{$\vcenter{\hbox{\tiny$\bullet$}}$}

\newcommand{\resumeSubHeadingListStart}{\begin{itemize}[leftmargin=0.15in, label={}]}
\newcommand{\resumeSubHeadingListEnd}{\end{itemize}}
\newcommand{\resumeItemListStart}{\begin{itemize}}
\newcommand{\resumeItemListEnd}{\end{itemize}\vspace{-5pt}}

%-------------------------------------------
%%%%%%  RESUME STARTS HERE  %%%%%%%%%%%%%%%%%%%%%%%%%%%%


\begin{document}

%----------HEADING----------
\begin{center}
    \textbf{\Huge \scshape Lorenzo Baladier} \\ \vspace{2pt}
    \small \href{mailto:contact@lorenzo-baladier.fr}{\underline{contact@lorenzo-baladier.fr}} $|$
    \href{https://linkedin.com/in/lorenzo-baladier}{\underline{linkedin.com/in/lorenzo-baladier}} $|$
    \href{https://github.com/Loren844}{\underline{github.com/Loren844}} $|$
    \href{https://lorenzo-baladier.fr}{\underline{lorenzo-baladier.fr}} \\ \vspace{2pt}
    \textbf{\small Ingénieur logiciel embarqué $\cdot$ Recherche un CDI à Paris à partir de septembre 2027}
\end{center}

%-----------FORMATION-----------
\section{Formation}
  \resumeSubHeadingListStart
    \resumeSubheading
      {Polytech Nice Sophia -- Université Côte d'Azur}{Sophia Antipolis, FR}
      {Diplôme d'ingénieur en informatique, option IoT \& Systèmes Cyber-Physiques (alternance)}{2024 -- 2027}
    \resumeSubheading
      {IUT Nice Côte d'Azur}{Sophia Antipolis, FR}
      {DUT Informatique (alternance)}{2022 -- 2024}
    \resumeSubheading
      {Lycée Pierre et Marie Curie}{Menton, FR}
      {Baccalauréat scientifique, section EsaBac (double diplôme franco-italien)}{2019 -- 2022}
  \resumeSubHeadingListEnd

%-----------EXPERIENCE-----------
\section{Expérience}
  \resumeSubHeadingListStart

    \resumeSubheading
      {Ingénieur logiciel embarqué \& SDV -- Alternant}{Sept. 2024 -- Aujourd'hui}
      {Ampere Software Technology}{Biot, FR}
      \resumeItemListStart
        \resumeItem{Conception et benchmark d'une plateforme de prototypage basée sur l'environnement \textbf{Software-Defined Vehicle} de Google : \textbf{CarOS}}
        \resumeItem{Bring-up de la plateforme sur \textbf{Raspberry Pi} et sur machines virtuelles \textbf{Cuttlefish}; mise en place des pipelines \textbf{CI/CD}}
        \resumeItem{Réalisation de cas d'usage de bout en bout : communication orientée services, diagnostic véhicule et traitement de données}
        \resumeItem{Réécriture d'un prototype Python en un outil \textbf{Rust} modulaire (CLI/TUI) générant des commandes de test \textbf{SOME/IP}, \textbf{Protobuf} à l'aide de GraphQL, testé par mocks à base de traits}
      \resumeItemListEnd

    \resumeSubheading
      {Stagiaire robotique \& asservissement visuel}{Mai 2026 -- Août 2026}
      {DNIIT -- Da Nang Institute of Innovative Technology}{Da Nang, Vietnam}
      \resumeItemListStart
        \resumeItem{Asservissement visuel en boucle fermée d'un bras \textbf{SCARA} industriel pour du pick-and-place de cibles en mouvement}
        \resumeItem{Intégration de la détection d'objets temps réel \textbf{YOLO} et de l'estimation de pose avec \textbf{ROS 2} et \textbf{MoveIt 2}}
        \resumeItem{Validation de la cinématique et de la commande PID sous \textbf{MATLAB} et \textbf{Gazebo} avant déploiement sur le robot réel -- \href{https://github.com/Loren844/project-10-dniit}{\underline{code source}}}
      \resumeItemListEnd

    \resumeSubheading
      {Développeur logiciel -- Alternant}{Sept. 2022 -- Juil. 2024}
      {ESI -- European Systems Integration}{Le Cannet, FR}
      \resumeItemListStart
        \resumeItem{Développement d'une application \textbf{Java} d'export JSON de données complexes depuis une base distante}
        \resumeItem{Automatisation des tests de bout en bout avec Selenium/TestNG, puis migration vers \textbf{TypeScript/Cypress}}
      \resumeItemListEnd

  \resumeSubHeadingListEnd

%-----------PROJETS-----------
\section{Projets}
    \resumeSubHeadingListStart
      \resumeProjectHeading
          {\textbf{Émulateur CHIP-8} $|$ \emph{Rust, SDL2} -- \href{https://github.com/Loren844/rust-chip-8}{\underline{github.com/Loren844/rust-chip-8}}}{2025}
          \resumeItemListStart
            \resumeItem{Émulation de la machine virtuelle CHIP-8 : boucle fetch-decode-execute, mémoire, registres, timers, affichage et clavier}
          \resumeItemListEnd
      \resumeProjectHeading
          {\textbf{Go-Biking} $|$ \emph{C\#, SOAP/REST, ActiveMQ}}{Projet d'école}
          \resumeItemListStart
            \resumeItem{Calculateur d'itinéraires orienté services combinant API de vélos en libre-service et de routage, avec proxy de cache et messagerie asynchrone}
          \resumeItemListEnd
    \resumeSubHeadingListEnd

%-----------COMPETENCES-----------
\section{Compétences}
 \begin{itemize}[leftmargin=0.15in, label={}]
    \small{\item{
     \textbf{Langages}{ : Rust, C, C++, Python, Java, TypeScript, Bash} \\
     \textbf{Embarqué \& automobile}{ : Linux, Raspberry Pi, Cuttlefish, Software-Defined Vehicle, SOA, SOME/IP, Protobuf} \\
     \textbf{Outils}{ : Git, GitLab CI/CD, Cypress, Selenium} \\
     \textbf{Langues}{ : Français (natif), Italien (courant, EsaBac), Anglais (TOEIC 955/990)}
    }}
 \end{itemize}

\end{document}
